% ---------------------------------------------------------------------------
% ---------------------------------------------------------------------------
% Junção de templates encontrados no Overleaf
% facilitando a vida do aluno de pós da UFABC
% Modelo LaTex para preparação do documento final de Dissertação de Mestrado
% ninguém do programa de pós da informação validou se presta
% ---------------------------------------------------------------------------
% ---------------------------------------------------------------------------

\documentclass[
  % -- opções da classe memoir --
  12pt, % tamanho da fonte
  openright, % capítulos começam em pág ímpar (insere página vazia caso preciso)
  twoside, % para impressão em verso e anverso. Oposto a oneside
  a4paper, % tamanho do papel.
  % -- opções da classe abntex2 --
  %chapter=TITLE,      % títulos de capítulos convertidos em letras maiúsculas
  %section=TITLE,      % títulos de seções convertidos em letras maiúsculas
  %subsection=TITLE,    % títulos de subseções convertidos em letras maiúsculas
  %subsubsection=TITLE,  % títulos de subsubseções convertidos em letras maiúsculas
  % -- opções do pacote babel --
  english, % idioma adicional para hifenização
  %french,          % idioma adicional para hifenização
  %spanish,        % idioma adicional para hifenização
  brazil % o último idioma é o principal do documento
]{abntex2}

% ---------------------
% Pacotes OBRIGATÓRIOS
% ---------------------
\usepackage{lmodern} % Usa a fonte Latin Modern
\usepackage[T1]{fontenc} % Selecao de codigos de fonte.
\usepackage[utf8]{
  inputenc
} % Codificacao do documento (conversão automática dos acentos)
\usepackage{lastpage} % Usado pela Ficha catalográfica
\usepackage{indentfirst} % Indenta o primeiro parágrafo de cada seção.
\usepackage{color} % Controle das cores
\usepackage{graphicx, graphicx} % Inclusão de gráficos
\usepackage{epsfig, subfig} % Inclusão de figuras
\usepackage{microtype} % Melhorias de justificação
% ---------------------

% ---------------------
% Pacotes ADICIONAIS
% ---------------------
\usepackage{lipsum} % Geração de dummy text
\usepackage{amsmath, amssymb, mathrsfs, mathabx} % Comandos matemáticos avançados
\usepackage{setspace} % Para permitir espaçamento simples, 1 1/2 e duplo
\usepackage{verbatim} % Para poder usar o ambiente "comment"
\usepackage{
  tabularx
} % Para poder ter tabelas com colunas de largura auto-ajustável
\usepackage{multirow} % Para mesclar linhas em tabelas
\usepackage{
  afterpage
} % Para executar um comando depois do fim da página corrente
\usepackage{url} % Para formatar URLs (endereços da Web)
% ---------------------

% ---------------------
% Pacotes de CITAÇÕES
% ---------------------
\usepackage[brazilian, hyperpageref]{backref} % Paginas com as citações na bibl
\usepackage[alf]{abntex2cite} % Citações padrão ABNT (alfa)
%\usepackage[num]{abntex2cite}        % Citações padrão ABNT (numericas)
% ---------------------

% O tamanho da identação do parágrafo é dado por:
\setlength{\parindent}{1.3cm}

% Controle do espaçamento entre um parágrafo e outro:
\setlength{\parskip}{0.2cm} % tente também \onelineskip

% Data da defesa
\data{Agosto de 2025}

% ---------------------
% Compila o indice
% ---------------------
\makeindex
% ---------------------

%%%%%%%%%%%%%%%%%%%%%%%%%%%
%%  INICIO DO DOCUMENTO  %%
%%%%%%%%%%%%%%%%%%%%%%%%%%%
\begin{document}
% Configurações de CITAÇÕES para abntex2
% ---
% CONFIGURAÇÕES DE PACOTES
% ---

% ---
% Configurações do pacote backref
% Usado sem a opção hyperpageref de backref
\renewcommand{\backrefpagesname}{Citado na(s) página(s):~}
% Texto padrão antes do número das páginas
\renewcommand{\backref}{}
% Define os textos da citação
\renewcommand*{\backrefalt}[4]{ \ifcase #1 %
  Nenhuma citação no texto.%
  \or Citado na página #2.%
  \else Citado #1 vezes nas páginas #2.%
\fi}%
% ---


% Inclusão de dados para CAPA e FOLHA DE ROSTO (título, autor, orientador, etc.)
% ---
% Informações de dados para CAPA e FOLHA DE ROSTO
% ---
\titulo{Comparativo de estratégias para lidar com desbalanceamento de dados na avaliação de risco de crédito}
\autor{William Fernandes Dias} \local{Santo André - SP} \orientador{Paulo Pisani}
\coorientador{} \instituicao{ Universidade Federal do ABC -- UFABC \par Centro de Matemática, Computação e Cognição \par Bacharelado em Ciência da Computação }
\tipotrabalho{Projeto de Graduação em Computação}
% O preambulo deve conter o tipo do trabalho, o objetivo,
% o nome da instituição e a área de concentração
\preambulo{\textbf{Projeto de Graduação em Computação} apresentado ao Centro de Ciências, Computação e Cognição como parte dos requisitos necessários para a obtenção do Título de Bacharel em Ciência da Computação.}
% ---

% Inclui Configurações de aparência do PDF Final
%  Configurações de aparência do PDF final

% alterando o aspecto da cor azul
\definecolor{blue}{RGB}{41,5,195}

% informações do PDF
\makeatletter
\hypersetup{
	%pagebackref=true,
	pdftitle={\@title},
	pdfauthor={\@author},
	pdfsubject={\imprimirpreambulo},
	pdfcreator={LaTeX with abnTeX2},
	pdfkeywords={abnt}{latex}{abntex}{abntex2}{trabalho acadêmico},
	linkcolor=blue, % color of internal links
	citecolor=blue, % color of links to bibliography
	filecolor=magenta, % color of file links
	urlcolor=blue,
	bookmarksdepth=4
}
\makeatother
% ---

% Retira espaço extra obsoleto entre as frases.
\frenchspacing

% ----------------------------------------------------------
% ELEMENTOS PRÉ-TEXTUAIS (Capa, Resumo, Abstract, etc.)
% ----------------------------------------------------------
\pretextual

% Capa
% ---
% Impressão da Capa
% ---
\begin{capa}%
  \begin{figure}[h!]%
    \centering%
    \includegraphics[scale=1.2]{figs/logo.png}%
  \end{figure}%
  \center \ABNTEXchapterfont\large{Universidade Federal do ABC
    \\ Centro de Matemática, Computação e Cognição \\ Bacharelado em
  Ciência da Computação}
  %\vspace{1.5cm}

  \vfill
  \ABNTEXchapterfont\bfseries\LARGE\imprimirtitulo
  \vfill

  %\vfill
  \ABNTEXchapterfont\large\imprimirautor
  \vfill

  \large\imprimirlocal, \large\imprimirdata

  \vspace*{1cm}
\end{capa}
% ---


% Folha de rosto (o * indica que haverá a ficha bibliográfica)
\imprimirfolhaderosto*

% Imprimir Ficha Catalografica
% % ---
% Ficha Catalográfica
% ---
% Isto é um exemplo de Ficha Catalográfica, ou ``Dados internacionais de
% catalogação-na-publicação''. Você pode utilizar este modelo como referência.
% Porém, talvez a biblioteca lhe fornece um PDF
% com a ficha catalográfica definitiva após a defesa do trabalho. Quando estiver
% com o documento, salve-o como PDF no diretório do seu projeto e substitua todo
% o conteúdo de implementação deste arquivo pelo comando abaixo:
%
% \begin{fichacatalografica}
%     \includepdf{fig_ficha_catalografica.pdf}
% \end{fichacatalografica}
\begin{fichacatalografica}
  \vspace*{\fill} % Posição vertical
  \hrule % Linha horizontal
  \begin{center} % Minipage Centralizado
    \begin{minipage}[c]{12.5cm} % Largura
      \imprimirautor

      \hspace{0.5cm} \imprimirtitulo / \imprimirautor. -- \imprimirlocal,
      \imprimirdata-

      \hspace{0.5cm} \pageref{LastPage} p. : il. (algumas color.) ; 30 cm.\\

      \hspace{0.5cm} \imprimirorientadorRotulo~\imprimirorientador\\

      \hspace{0.5cm}
      \parbox[t]{\textwidth}{\imprimirtipotrabalho~--~\imprimirinstituicao,
      \imprimirdata.}
      \\

      \hspace{0.5cm} 1. Palavra-chave1. 2. Palavra-chave2. I. Orientador. II.
      Universidade xxx. III. Faculdade de xxx. IV. Título\\

      \hspace{8.75cm} CDU 02:141:005.7\\
    \end{minipage}
  \end{center}
  \hrule
\end{fichacatalografica}
% ---


% Inserir Folha de Aprovação
% % ---
% Assinaturas
% ---
% Isto é um exemplo de Folha de aprovação, elemento obrigatório da NBR
% 14724/2011 (seção 4.2.1.3). Você pode utilizar este modelo até a aprovação
% do trabalho. Após isso, substitua todo o conteúdo deste arquivo por uma
% imagem da página assinada pela banca com o comando abaixo:
%
% \includepdf{folhadeaprovacao_final.pdf}
%
\begin{folhadeaprovacao}
    \begin{center}
        {\ABNTEXchapterfont\large\imprimirautor}

        \vspace*{\fill}
        \vspace*{\fill}
        \begin{center}
            \ABNTEXchapterfont\bfseries\Large\imprimirtitulo
        \end{center}
        \vspace*{\fill}

        \hspace{.45\textwidth}
        \begin{minipage}{.5\textwidth}
            \imprimirpreambulo
        \end{minipage}%
        \vspace*{\fill}
    \end{center}

    % Isso na versao final do trabalho!!!
    Trabalho aprovado. \imprimirlocal, 01 de janeiro de 2014:

    \assinatura{\textbf{\imprimirorientador} \\ Orientador} \assinatura{\textbf{\imprimircoorientador} \\ Co-Orientador}
    \assinatura{\textbf{Professor} \\ Convidado 1} \assinatura{\textbf{Professor} \\ Convidado 2}
    \assinatura{\textbf{Professor} \\ Convidado 3}

    \begin{center}
        \vspace*{0.5cm}
        {\large\imprimirlocal} \par {\large\imprimirdata}
        \vspace*{1cm}
    \end{center}
\end{folhadeaprovacao}
% ---

% Dedicatória
% % ---
% Dedicatória
% ---
\begin{dedicatoria}
  \vspace*{\fill}
  \centering
  \noindent
  \textit{ Aos verme que roeu as frias carnes de meu cadáver.}
  \vspace*{\fill}
\end{dedicatoria}
% ---


% Agradecimentos
% % ---
% Agradecimentos
% ---
\begin{agradecimentos}
    Agradeço a Xuxa, meus pais, cachorro, gato e papagaio, por ...

    Agradeço ao meu orientador, XXXXXXXXX, por todos os conselhos, pela paciência
    e ajuda nesse período.

    Aos meus amigos ...

    Aos professores ...

    À XXXXXX pelo apoio financeiro para realização deste trabalho de pesquisa.
\end{agradecimentos}
%% ---

% Epígrafe
% % ---
% Epígrafe
% ---
\begin{epigrafe}
  \vspace*{\fill}
  \begin{flushright}
    \textit{``Não sei o que, \\ não sei o que,\\ não sei o que lá.''\\ (Autor
    Desconhecido)}
  \end{flushright}
\end{epigrafe}
% ---


% Resumo e Abstract
% ---
% RESUMOS
% ---

% RESUMO em português
\setlength{\absparsep}{18pt} % ajusta o espaçamento dos parágrafos do resumo
\begin{resumo}
  A avaliação de risco de crédito é essencial para instituições financeiras,
  permitindo a concessão responsável de empréstimos e a redução de
  prejuízos. Com
  o avanço do aprendizado de máquina, modelos preditivos tornaram-se
  ferramentas valiosas para identificar possíveis casos de
  inadimplência, auxiliando
  na tomada de decisões mais precisas e reduzindo drasticamente o
  trabalho manual.
  No entanto, a previsão de inadimplência é desafiada pelo
  desequilíbrio dos conjuntos
  de dados, onde casos de calote são significativamente mais raros do que
  pagamentos bem-sucedidos. Esse desbalanceamento leva a modelos tendenciosos,
  que frequentemente falham em identificar tomadores de alto risco com
  precisão. Para mitigar esse problema, técnicas como re-amostragem de dados e
  aprendizagem sensível ao custo são amplamente utilizadas. Este projeto tem o
  objetivo de realizar uma avaliação comparativa de técnicas para
  lidar com desbalanceamento
  de dados no contexto de avaliação de risco de crédito.

  \textbf{Palavras-chaves}: Previsão de inadimplência. Desbalanceamento de
  dados. Aprendizado sensível ao custo. Classes desbalanceadas. Risco
  de crédito.
\end{resumo}

% ABSTRACT in english
\begin{resumo}
  [Abstract]
  \begin{otherlanguage*}
    {english} Credit risk assessment is essential for financial institutions,
    enabling responsible lending and loss reduction. With the advancement of
    machine learning, predictive models have become valuable tools for
    identifying potential default cases, supporting more accurate
    decision-making
    and drastically reducing manual work. However, default prediction
    is challenged
    by the imbalance in datasets, where default cases are significantly
    rarer than successful payments. This imbalance leads to biased models that
    often fail to accurately identify high-risk borrowers. To mitigate this
    issue, techniques such as data resampling and cost-sensitive learning are
    widely used. This project aims to perform a comparative evaluation of
    techniques for handling data imbalance in the context of credit risk
    assessment.

    \vspace{\onelineskip}

    \noindent
    \textbf{Keywords}: Default prediction. Data imbalance.
    Cost-sensitive learning.
    Imbalanced classes. Credit risk.
  \end{otherlanguage*}
\end{resumo}


% Lista de ilustrações
\pdfbookmark[0]{\listfigurename}{lof} \listoffigures*
\cleardoublepage

% Lista de tabelas
\pdfbookmark[0]{\listtablename}{lot} \listoftables*
\cleardoublepage

% Inserir o SUMÁRIO
\pdfbookmark[0]{\contentsname}{toc}
\tableofcontents
*
\cleardoublepage

% ----------------------------------------------------------
% ELEMENTOS TEXTUAIS (Capítulos)
% ----------------------------------------------------------
\textual
% Elementos textuais com numeração arábica
\pagenumbering{arabic}
% Reinicia a contagem do número de páginas
\setcounter{page}{1}

% Inclui cada capitulo da Dissertação
% ----------------------------------------------------------
% Introdução
% Capítulo sem numeração, mas presente no Sumário
% ----------------------------------------------------------

\chapter[Introdução]{Introdução}
\addcontentsline{toc}{chapter}{Introdução}

% Este documento segue as normas estabelecidas
% pela~\citeonline[3.1-3.2]{NBR6028:2003}.

% \section*{Figuras}\label{sec:figuras}
% \addcontentsline{toc}{section}{figuras}

% As normas da~\citeonline[3.1-3.2]{NBR6028:2003} especificam que o
% caption da figura deve vir abaixo da mesma.

% A Figura~\ref{fig:log} ilustra...

% \begin{figure}[htpb]
%    \centering
%    \includegraphics[scale=.3]{figs/logo}
%    \caption{Breve explicação sobre a figura. Deve vir abaixo da mesma.}
%    \label{fig:log}
% \end{figure}

% \section*{Tabelas}\label{sec:tabelas}
% \addcontentsline{toc}{section}{tabelas}

% A Tabela~\ref{tab:tabela} apresenta os resultados...

% \begin{table}[htpb]
%    \centering
%    \caption{Breve explicação sobre a tabela. Deve vir acima da
% mesma.}\label{tab:tabela}
%    \begin{tabular}{|l|c|c|c|c|c|c|r|}
%         \hline
%         \small{XX} & \small{FF} & \small{PP} & \small{YY} &
% \small{Yr} & \small{xY} & \small{Yx} & \small{ZZ} \\ \hline
%                615 &    18      &     2558   &    0,9930  &
% 0,9930  &    0,9930  &    0,9930  &    0,9930  \\ \hline
%                615 &    18      &     2558   &    0,9930  &
% 0,9930  &    0,9930  &    0,9930  &    0,9930  \\ \hline
%                615 &    18      &     2558   &    0,9930  &
% 0,9930  &    0,9930  &    0,9930  &    0,9930  \\ \hline
%                615 &    18      &     2558   &    0,9930  &
% 0,9930  &    0,9930  &    0,9930  &    0,9930  \\ \hline
%                615 &    18      &     2558   &    0,9930  &
% 0,9930  &    0,9930  &    0,9930  &    0,9930  \\ \hline
%    \end{tabular}
% \end{table}

A avaliação de risco de crédito é um componente essencial para a estabilidade do mercado financeiro, pois permite que instituições avaliem com precisão a probabilidade de inadimplência de clientes. Erros nessa avaliação podem resultar não apenas em prejuízos financeiros individuais, mas também em crises sistêmicas, como evidenciado na crise do subprime de 2008.

Nos últimos anos, o uso de algoritmos de aprendizado de máquina tem se consolidado como uma alternativa promissora aos modelos estatísticos tradicionais, como a análise discriminante linear e a regressão logística. Tais algoritmos oferecem maior flexibilidade na modelagem de dados complexos e têm apresentado desempenho superior em diversos cenários de previsão de crédito~\cite{Shi2022}.

No entanto, um desafio persistente nesse domínio é o desbalanceamento das classes: casos de inadimplência são, geralmente, muito menos frequentes do que pagamentos em dia. Esse desequilíbrio tende a enviesar os modelos preditivos, reduzindo sua capacidade de identificar tomadores de alto risco, justamente os casos mais relevantes na prática~\cite{Namvar2018}.

Para lidar com esse problema, diversas estratégias têm sido propostas. Dentre elas, destacam-se as técnicas de re-amostragem, como o SMOTE e o \textit{under-sampling}, que tentam ajustar a distribuição das classes no conjunto de dados; e abordagens de aprendizado sensível ao custo, como o \textit{MetaCost}, que buscam incorporar os impactos assimétricos dos erros de classificação ao próprio processo de aprendizagem~\cite{FernndezCs2018,Wei2025}.

Este trabalho tem como objetivo realizar um estudo comparativo entre diferentes estratégias para lidar com o desbalanceamento de dados na avaliação de risco de crédito. Serão considerados cenários variados, incluindo crédito pessoal, financiamento imobiliário e crédito corporativo, a fim de avaliar a robustez das abordagens analisadas em contextos distintos. Ao final, espera-se oferecer subsídios para a escolha de técnicas mais adequadas à construção de modelos preditivos eficazes em situações de classes desbalanceadas.

\section{Justificativa}\label{sec:motivacao}
\addcontentsline{toc}{section}{Motivação}

\section{Objetivos}\label{sec:objetivos}
\addcontentsline{toc}{section}{Objetivos}


% PARTE - Define a divisão do documento em partes (Não é obrigatório)
% \part{Fundamentação Teórica}
\include{capitulos/referencias}
\chapter{Fundamentação teórica}\label{cap:fundamentacao-teorica}

No setor financeiro, a avaliação de risco de crédito desempenha um papel fundamental na manutenção da estabilidade, rentabilidade e confiança nas operações~\cite{Chakraborty2020}. Esse processo consiste em avaliar a probabilidade de um tomador de empréstimo não cumprir com suas obrigações, sendo crucial para a tomada de decisões de crédito informadas e para a gestão do risco geral da carteira de uma instituição. Tradicionalmente, as instituições financeiras se baseiam em uma variedade de métodos para mensurar o risco, cada um com suas vantagens e desvantagens~\cite{Bello2023}. Entre as abordagens mais consolidadas, destacam-se:

\begin{itemize}
  \item \textbf{Modelos de pontuação de crédito (\textit{credit scoring}):} Técnicas estatísticas que atribuem um valor numérico a cada tomador para indicar a probabilidade de inadimplência~\cite{Pang2021}. Essa pontuação é calculada com base em fatores como histórico de pagamentos, total de crédito utilizado em relação ao limite disponível, tempo de relacionamento com o mercado de crédito, diversidade de tipos de crédito e consultas abrangentes ao seu histórico.

  \item \textbf{Análise de indicadores financeiros:} Métricas quantitativas extraídas das demonstrações contábeis de uma empresa para avaliar seu desempenho e solvência. Esses indicadores incluem rácios de alavancagem, como a relação dívida/patrimônio líquido~\cite{Coulon2019}; de lucratividade, como o retorno sobre os ativos; e de liquidez, como o índice de liquidez corrente.

  \item \textbf{Métodos estatísticos:} Técnicas como a regressão logística e a análise discriminante são utilizadas para modelar a relação entre as características de um indivíduo ou empresa e o risco de crédito associado.
\end{itemize}

Contudo, essas abordagens tradicionais apresentam limitações significativas. Elas podem ser muito subjetivas e suscetíveis a vieses, visto que dependem da interpretação humana dos dados~\cite{Bello2023, Goel2021}. Além disso, o uso de dados históricos pode não ser suficiente para antecipar tendências em mercados que estão em constante evolução. A inflexibilidade para incorporar novas informações e a dificuldade de se adaptar a novas condições de mercado comprometem seu poder preditivo. Essa falta de adaptabilidade, somada a um escopo frequentemente limitado a fatores quantitativos, tem posicionado as técnicas baseadas em aprendizado de máquina como uma alternativa poderosa no campo da análise de risco de crédito~\cite{Mhlanga2021}.

Aprendizado de máquina é o processo de fazer com que computadores modifiquem, ou adaptem, as suas ações de modo que elas se tornem progressivamente mais eficazes segundo um conjunto de métricas~\cite{StephenMarsland2014}. É uma sub-área da Inteligência Artificial que foca na reprodução computacional do aprendizado, no sentido de desenvolver programas capazes de generalizar um problema: reconhecer que dentro de um determinado contexto (entrada), se uma decisão em particular (saída) foi a correta, então ela pode funcionar novamente, ou se ela foi equivocada, uma estratégia diferente deve ser elaborada.

Os métodos utilizados para atingir esse comportamento podem ser divididos em aprendizado supervisionado, não-supervisionado e por reforço. No aprendizado supervisionado, um conjunto de dados de treinamento é utilizado juntamente com as respostas esperadas (rótulos) é fornecido, e com isso um algoritmo generaliza o problema, se tornando capaz de prever a resposta correta para novas entradas. Quando essa resposta é categórica, diz-se que o algoritmo é de classificação, e quando é contínua chama-se regressão~\cite{SindhuMeena2020}. O aprendizado não-supervisionado, por sua vez, opera em dados não rotulados para identificar padrões~\cite{Dike2018}. No aprendizado por reforço, um agente autônomo aprende a atingir um objetivo por meio de tentativa e erro enquanto interage com um ambiente: o agente utiliza o feedback que recebe na forma de recompensas para aprender uma política de ações que maximize a recompensa total a longo prazo~\cite{Sutton2018}.

O escopo deste projeto concentra-se na tarefa de classificação binária aplicada a conjuntos de dados utilizados na avaliação de risco de crédito. Nesse contexto, é comum a ocorrência de desbalanceamento entre as classes, o que pode impactar negativamente o desempenho dos modelos preditivos~\cite{Namvar2018}. Assim, estratégias que considerem esse desbalanceamento tendem a produzir resultados mais robustos. Este capítulo, portanto, explora os principais conceitos e técnicas relacionados ao desenvolvimento de modelos de classificação em cenários com dados desbalanceados.

\section{Métricas de desempenho de classificadores}

A escolha das métricas de desempenho é determinante para uma avaliação precisa e confiável de um classificador~\cite{Gaudreault2021}. Os indicadores de desempenho preditivo utilizados hoje têm sua origem em variados domínios incluindo estatística (distância, similaridade binária etc.), processamento de sinais (\textit{area under the receiver operating characteristic curve}) AUC, recuperação de informação (\textit{precision}, \textit{recall}), medicina diagnóstica (sensibilidade, especificidade), reconhecimento estatístico de padrões (acurácia), entre outros~\cite{Canbek2023}. Quando se trata da avaliação de classificadores binários, as métricas utilizadas na literatura podem ser agrupadas segundo a sua interpretação do erro~\cite{Ferri2009}:  acurácia e \textit{F1-score}, por exemplo, representam o erro de forma qualitativa utilizando proporções entre a contagem de erros e acertos, enquanto que média do erro absoluto e \textit{Brier Score} dão uma visão probabilística do erro, isto é, medem o desvio das previsões da sua probabilidade real.

Para o uso das métricas qualitativas, a contagem dos erros e acertos é representada com uma matriz \(M\), chamada de Matriz de confusão. Nesta matriz, cada elemento \(M_{ij}\) corresponde ao número de vezes com que uma instância da classe \(i\) foram identificados como pertencentes à classe \(j\).

Em problemas de classificação binária desbalanceada, a classe mais frequente é chamada de majoritária, e a classe menos frequente é chamada de minoritária~\cite{Seiffert2008,Batuwita2010}. Com base nisso, a matriz de confusão para esses casos pode ser vista na Tabela~\ref{tab:matriz-confusao}.

\begin{table}[htbp]
  \centering
  \begin{tabular}{cc|cc|}
    \cline{3-4}
    \multicolumn{2}{c|}{\multirow{2}{*}{}}                   & \multicolumn{2}{c|}{\textbf{Real}}                                                                                \\ \cline{3-4}
    \multicolumn{2}{c|}{}                                    & \multicolumn{1}{c|}{\textbf{Negativa}} & \textbf{Positiva}                                                        \\ \hline
    \multicolumn{1}{|c|}{\multirow{2}{*}{\textbf{Predição}}} & \textbf{Negativa}                      & \multicolumn{1}{c|}{TN (verdadeiro negativo)} & FN (falso negativo)      \\ \cline{2-4}
    \multicolumn{1}{|c|}{}                                   & \textbf{Positiva}                      & \multicolumn{1}{c|}{FP (falso positivo)}      & TP (verdadeiro positivo) \\ \hline
  \end{tabular}
  \caption{Matriz de confusão binária.}
  \label{tab:matriz-confusao}
\end{table}

\section{Técnicas de re-amostragem}

Considerando a perda de desempenho que muitos algoritmos apresentam diante de dados desbalanceados, alguns métodos foram desenvolvidos para mitigar a discrepância entre as classes e evitar o surgimento de viés a favor da classe majoritária. Esse processo, conhecido de re-amostragem, altera o conjunto de dados adicionando ou removendo instâncias seguindo algum critério~\cite{Chakravarthy2019}. Essas abordagens são amplamente aplicadas devido a sua versatilidade: como envolvem apenas os dados, o seu uso independe de qualquer modelo específico ou contexto de classificação~\cite{Carvalho2025}. As técnicas de re-amostragem podem ser divididas em três tipos~\cite{Haixiang2017}:

\begin{itemize}
  \item sobre-amostragem, (ou \textit{oversampling}) consiste em criar novas instâncias da classe minoritária, reduzindo o grau de desbalanceamento do conjunto de dados;
  \item sub-amostragem (ou \textit{under-sampling}), que remove amostras da classe majoritária~\cite{Mohammed2020};
  \item abordagem híbrida, que combina os dois tipos anteriores.
\end{itemize}

Dentre as técnicas mais simples — e ainda amplamente utilizadas — destacam-se as não heurísticas, como o \textit{Random Oversampling} (ROC), que gera amostras sintéticas, e o \textit{Random Undersampling}, que remove aleatoriamente instâncias do conjunto de dados durante o treinamento. No entanto, esses métodos não garantem uma melhoria de desempenho em todas as aplicações~\cite{Yang2024} e podem acarretar problemas, como \textit{overfitting} ou perda de informações relevantes~\cite{Fernndez2018}.

Partindo para o campo das técnicas heurísticas temos, por exemplo, o \textit{Tomek links}, um método de sub-amostragem que busca melhorar a performance de classificadores reduzindo ambiguidades em pontos de borda do conjunto de dados. Sejam \(S_{i}\) e \(S_{j}\) dois exemplos de classes distintas, e seja \(d_{ij}\) a distância Euclidiana entre eles. Um par \(S_{i}, S_{j}\) é considerado um \textit{Tomek link} se não existir nenhum \(S_{l}\) tal que \(d_{il} < d_{ij}\) ou \(d_{jl} < d_{ij}\)~\cite{Batista2004}. Para aplicar o \textit{Tomek links}, primeiro se calcula os vizinhos mais próximos de cada membro da classe minoritária. Em seguida, identificam-se os pares. Para cada par, o elemento da classe majoritária é eliminado. Esse processo pode ser repetido iterativamente para aumentar a separação entre as classes~\cite{Arichandrapandian2024}. Apesar de ter sido proposto inicialmente para problemas de classificação de rótulo único, ele já foi adaptado para lidar com múltiplas classes~\cite{Pereira2020}.

\textit{Synthetic minority oversampling} (SMOTE), é um método de sobre-amostragem que introduz uma abordagem baseada em interpolação para gerar novas instâncias da classe minoritária~\cite{Wei2025}. Ao invés de criar exemplos totalmente aleatórios, o SMOTE cria amostras sintéticas a partir da combinação de instâncias próximas, o que contribui para uma modelagem mais robusta e representativa da classe minoritária, além de melhorar a capacidade de generalização dos modelos preditivos. Inicialmente, o método agrupa as amostras positivas utilizando, por exemplo, o algoritmo \textit{k-Nearest Neighbors} (k-NN). Em seguida, para cada nova amostra a ser gerada, duas instâncias \(x_{1}\) e \(x_{2}\) pertencentes ao mesmo grupo são selecionadas, e a instância sintética é obtida por meio da interpolação \(x = x_{1} + \lambda \bigtimes (x_{1} - x_{2})\), onde \(\lambda\) é um número aleatório no intervalo \([0, 1]\). Essa técnica tem se mostrado eficaz na redução do viés de aprendizado causado pelo desbalanceamento, é amplamente utilizada em diferentes domínios de aplicação e inspirou diversos outros algoritmos de re-amostragem para lidar com desbalanceamento~\cite{Fernndez2018}. \textit{Adaptative synthetic} (ADASYN) \textit{sampling} é uma alternativa ao SMOTE para sobre-amostragem. Nessa abordagem, a distribuição das instâncias negativas é utilizada para se definir o número ideal de amostras sintéticas a serem geradas para cada amostra da classe minoritária com o objetivo de reduzir o viés introduzido pelo desbalanceamento, e ao mesmo tempo otimizar a precisão das decisões~\cite{He2008}.

Em bases de dados desbalanceadas, é comum que aglomerados da classe minoritária não sejam bem definidos devido a presença de elementos da classe majoritária. O contrário, onde instâncias negativas se encontram isoladas em uma parte do espaço amostral dominada por instâncias negativas, também é frequente na prática, inclusive após o uso de sobre-amostragem. O risco de \textit{over-fitting} aumenta em ambos os casos~\cite{Batista2004}. Para se atingir uma separação mais bem-definida entre as classes, uma combinação de técnicas de sobre-amostragem e sub-amostragem podem ser aplicadas, com a primeira servindo para aumentar a representatividade da classe minoritária, e a segunda para limpar o conjunto de dados, removendo redundâncias e reforçando fronteiras~\cite{Pereira2020}. SMOTE-Tomek e SMOTE-ENN são exemplos de técnicas híbridas, que adicionam uma fase de pós-processamento ao SMOTE, utilizando respectivamente \textit{Tomek links} e \textit{Edited Nearest Neighbor} (ENN).

\section{Aprendizado sensível ao custo}

O aprendizado sensível a custos é um conjunto de algoritmos que introduz pesos diferenciados para os erros em diferentes estágios da fase de treinamento~\cite{FernndezCs2018}. Os custos podem ser atribuídos aos atributos, quando o objetivo é otimizar a performance da classificação utilizando um subconjunto específico de atributos~\cite{Zhou2016}, ou às classes em problemas de classificação, em que falhas na identificação de instâncias da classe de interesse são penalizadas de forma mais severa durante o processo de treinamento.

No contexto de desbalanceamento, o aprendizado sensível a custos é uma abordagem a nível de algoritmo~\cite{Krawczyk2016}. Ela se baseia no fato de que a função de perda padrão, que atribui 0 para acertos e 1 para erros, é problemática em distribuições desequilibradas, pois é facilmente minimizada concentrando-se apenas na classe majoritária, negligenciando a classe minoritária ou até mesmo ignorando ela por completo, a depender do grau de desbalanceamento~\cite{Landgrebe2004}. Nessa abordagem, esse problema é resolvido adaptando-se a função de perda usando, por exemplo, uma matriz de custos \(C\), que associa um peso específico para cada \(C_{ij}\) possível. Com penalidades maiores para erros envolvendo a classe minoritária, o processo de treinamento dá uma enfase maior nessas instâncias. Comparada com técnicas de re-amostragem, o aprendizado sensível ao custo apresenta algumas vantagens, como a preservação do conjunto de dados original, capacidade de associar custos distintos para qualquer combinação de erros possível e eficiência computacional, o que a torna mais adequada para aplicações de \textit{Big Data}~\cite{Haixiang2017}.

Naturalmente, a efetividade do aprendizado sensível a custos depende fortemente da matriz de custos fornecida~\cite{FernndezCs2018}. Geralmente, ela é definida com o acompanhamento de um especialista no problema estudado. Sahin, Bulkan e Dulman (\citeyear{Sahin2013}), por exemplo, elaboraram uma matriz de custos para análise de risco de crédito após entrevistarem colaboradores de uma instituição financeira. Quando não é possível consultar um especialista, alguns trabalhos utilizam como alternativa estratégias analíticas para a definição de \(C\). Bahnsen, Djamila e Ottersten (\citeyear{Bahnsen2015}) desenvolveram um método para determinar o custo de cada exemplo durante a fase de treinamento.

Os algoritmos de aprendizado sensível ao custo podem ser categorizados entre aprendizagem direta, em que algoritmos existentes são adaptados para o uso de uma função de custo específica, e meta-aprendizado, que consiste em incorporar modificar a entrada (dados de teste) ou saída de um classificador~\cite{FernndezCs2018}. \textit{Cost-sensitive tree} é um exemplo de algoritmo da primeira categoria, que permite a introdução de custos, tanto dos atributos quanto de classificações incorretas, a uma árvore de decisão~\cite{Ling2006}.

Proposto por Domingos (\citeyear{Domingos1999}), \textit{MetaCost} é um método genérico de meta-aprendizado que permite adicionar um procedimento de minimização de custo a qualquer classificador, independentemente do formato de sua saída. A flexibilidade e a simplicidade da sua implementação, faz com o algoritmo se destaque entre os demais~\cite[text]{FernndezCs2018, Wang2018}. O processo de treinamento utilizando o \textit{MetaCost} é divido em duas partes. Primeiro, \(m\) subconjuntos são criados a partir do conjunto de dados de teste original, cada qual usado para treinar um classificador base (por exemplo, regressão logística, \textit{decision tree} etc.). Cada classificador gera probabilidades de classe para cada instância, ou apenas o rótulo quando ele não suporta probabilidades~\cite{Araf2024}. Em seguida, para cada exemplo original, é calculada a média das previsões dos classificadores base (ou votação por maioria), que é usada para calcular o custo esperado de classificá-lo em cada classe, multiplicando pela matriz de custos. O rótulo da instância é então redefinido para a classe minimiza o custo. Com o conjunto de treinamento re-rotulado, um classificador final é então treinado com os exemplos atualizados. Esse classificador final é sensível ao custo e tem um desempenho superior aos classificadores treinados com os dados originais~\cite{Chen2021}.

\section{Trabalhos Relacionados}

Ao longo dos anos, com o avanço acelerado das técnicas de modelagem, da capacidade computacional e dos softwares que sustentam aplicações de inteligência artificial, o uso de modelos baseados em \textit{machine learning} para análise de risco de crédito tem se tornado cada vez mais prevalente em relação às abordagens puramente estatísticas. Adicionalmente, o surgimento de novas modalidades de crédito, como o empréstimo entre pessoas (\textit{peer-to-peer} ou P2P), aliado ao uso de estratégias de diversificação cada vez mais sofisticadas, demanda modelos mais robustos, confiáveis e capazes de lidar com diferentes características dos dados.

Dado que a inadimplência é geralmente um evento raro, o desequilíbrio entre classes continua sendo um dos principais desafios na modelagem preditiva. Dessa forma, a experimentação e comparação de técnicas que lidam com o desbalanceamento ainda são de grande relevância. Esta seção apresenta uma revisão de estudos recentes que abordam a classificação de risco de crédito em cenários de forte desbalanceamento entre classes.

Namvar \textit{et al.}~(\citeyear{Namvar2018}) exploram esse problema no contexto de empréstimos P2P, comparando diferentes técnicas de reamostragem quanto à sua eficácia na melhoria do desempenho dos classificadores, com foco principal na métrica \textit{G-mean}. Os autores destacam o uso de \textit{Random Forest} combinado com \textit{Random Undersampling} (RUS) como uma solução promissora para a previsão de inadimplência. Apesar da redução na acurácia média, os modelos ajustados para o desbalanceamento mostraram menor viés em relação à classe majoritária, tornando-se mais adequados para o objetivo da tarefa.

Wang, Kou e Peng~(\citeyear{Namvar2018}), utilizando o mesmo conjunto de dados, avaliaram o impacto do aprendizado sensível ao custo por meio da técnica \textit{MetaCost}. Embora com foco em uma abordagem distinta, os resultados convergem com os de Namvar \textit{et al.}, reforçando que o tratamento do desbalanceamento durante o treinamento contribui para a redução do viés dos modelos, promovendo classificações mais equilibradas.

No âmbito do crédito corporativo, o problema do desbalanceamento se intensifica devido à maior complexidade dos dados e à necessidade de considerar um volume substancial de variáveis. Wei~(\citeyear{Wei2025}) propõe uma extensão do SMOTE, adaptada a esse cenário, na qual são atribuídos pesos diferenciados aos atributos mais relevantes para avaliação de risco em setores específicos. Apesar dos resultados promissores obtidos em múltiplos conjuntos de dados, o autor destaca que ainda existem relações complexas e intrínsecas entre as variáveis — como a dinâmica competitiva entre empresas e particularidades setoriais — que as abordagens atuais não conseguem generalizar adequadamente. Como alternativa, o uso de técnicas de \textit{Deep Learning} é apontado como uma possível direção para superar essas limitações.

Diante desse panorama, este trabalho tem como objetivo apresentar uma comparação sistemática entre as principais abordagens utilizadas para lidar com o desbalanceamento entre classes no contexto da análise de risco de crédito. Serão avaliadas diferentes técnicas em conjuntos de dados provenientes de distintas modalidades de empréstimo, com o intuito de identificar quais métodos apresentam melhor desempenho conforme métricas adequadas à classificação de risco.

\chapter{Materiais e Métodos}
\label{cap:ferramentas}

% \lipsum[43-45]

\section{Considerações Finais}

% \lipsum[23]

% PARTE
% \part{Resultados}
\chapter{Resultados e Discussão}
\label{cap:resultados}

% \lipsum[73]

\section{Base de Dados}

% \lipsum[72]

\section{Considerações Finais}

% \lipsum[74]

\chapter{Conclusões e Trabalhos Futuros}
\label{cap:conclusao}
\addcontentsline{toc}{chapter}{Conclusão e Trabalhos Futuros}

% \lipsum[81]

\section{Conclusões}

% \lipsum[82-84]

\section{Trabalhos Futuros}

% \lipsum[85]

% ----------------------------------------------------------
% ELEMENTOS PÓS-TEXTUAIS (Referências, Glossário, Apêndices)
% ----------------------------------------------------------
\postextual

% Referências bibliográficas
\bibliography{bibliografia}

% Glossário (Consulte o manual)
%\glossary

% Apêndices
% % ----------------------------------------------------------
% Apêndices
% ----------------------------------------------------------

% ---
% Inicia os apêndices
% ---
\begin{apendicesenv}
    % Imprime uma página indicando o início dos apêndices
    \renewcommand{\partapendices}{\part*{Apêndices}}
    \partapendices

    % ----------------------------------------------------------
    \chapter{Primeiro Apêncice}
    % ----------------------------------------------------------

    \lipsum[50] % Texto qualquer. REMOVER!!

    % ----------------------------------------------------------
    \chapter{Segundo apêndice com título tão grande quanto se queira porque ele
    já faz a quebra de linha da coisa toda}
    % ----------------------------------------------------------
    \lipsum[51-53] % Texto qualquer. REMOVER!!
\end{apendicesenv}
% ---

% Anexos
% % ----------------------------------------------------------
% Apêndices
% ----------------------------------------------------------

% ---
% Inicia os anexos
% ---
\begin{anexosenv}
  % Imprime uma página indicando o início dos anexos
  \renewcommand{\partanexos}{
  \part*{Anexos}}
  \partanexos

  % ---
  \chapter{Nome do Primeiro Anexo}
  % ---
  \lipsum[30] % Texto qualquer. REMOVER!!

  % ---
  \chapter{Nome de Outro Anexo}
  % ---

  \lipsum[32] % Texto qualquer. REMOVER!!
\end{anexosenv}


% Índice remissivo (Consultar manual)
%\printindex
\end{document}